\section{Recorded system results and development interventions}
\label{sec:recorded-results}
These tables preserve earlier system configurations and selected development
interventions. These seven tables are restored to the main text; this file preserves their
previous appendix organization for reference. Program preparation in the existing
composition cases is conditional on supplied public snippets and researcher
fixture wrappers, not an evaluated automatic issue-to-program pipeline.

\subsection{Experiment 1: repair effectiveness}
\label{sec:recorded-effectiveness}
The latest reported GPT-5.4/mini-SWE-agent comparison resolves 346/500 tasks
with ContextGraph versus 309/500 without memory: a 7.4 percentage-point gain
(Table~\ref{tab:verified500}). It resolves 19 more tasks than FAISS and 26 more
than Agent KB. These results motivate using experience, but do not isolate
which part of executable memory produces the gain. We therefore turn to
controlled development interventions, then return to cross-model and
efficiency results in Section~\ref{sec:recorded-generalization}.

\begin{table}[!htbp]
\centering
\caption{Recorded SWE-bench Verified500 results with GPT-5.4 and mini-SWE-agent.
These precede ContextGraph. Rates use 500 tasks; gains are points over no memory.}
\label{tab:verified500}
\begin{tabular}{lrrr}
\toprule
Method & Resolved & Rate & Gain (pp) \\
\midrule
No memory & 309 & 61.8\% & --- \\
ExpeL & 312 & 62.4\% & +0.6 \\
Agent KB & 320 & 64.0\% & +2.2 \\
FAISS & 327 & 65.4\% & +3.6 \\
ContextGraph & \textbf{346} & \textbf{69.2\%} & \textbf{+7.4} \\
\bottomrule
\end{tabular}
\end{table}

\subsection{Experiment 2: memory content and graph mechanisms}
\label{sec:recorded-interventions}
\paragraph{History supplies an omitted behavior.}
Table~\ref{tab:historicalrequirements} summarizes three complete-agent comparisons.
Django14404 receives the first-ranked source issue and patch; Django12708 and
Django11138 come from a fixed six-task cohort receiving three unchanged Flat
rules per task. Both arms pass official verification on all six cohort tasks.
The displayed pairs differ on historical requirements: preserving query strings,
removing redundant uniqueness, and honoring an explicit destination timezone.
These are one-repair-per-arm development comparisons, not a benchmark-wide
estimate of Joint success.

\begin{table}[!htbp]
\centering
\caption{Joint current/historical behavior in complete-agent comparisons.
Each displayed repair passes official verification. The columns compare the actual repairs with and without memory; the rows retain their distinct task requirements.}
\label{tab:historicalrequirements}
\begin{tabular}{llcc}
\toprule
Target & Historical addition & No memory & Memory \\
\midrule
Django14404 & Preserve query string & Fail & Pass \\
Django12708 & Delete redundant uniqueness & Fail & Pass \\
Django11138 & Honor destination timezone & Fail & Pass \\
\bottomrule
\end{tabular}
\end{table}

The additions cover three ways in which history can matter. The redirect
example concerns the same API; the uniqueness example constrains a neighboring
operation that must preserve a primary key and its data; the timezone example
constrains an independent argument of the conversion. In the last case, two
fresh agents refine the same completed patch. Additional issue-based review
leaves it unchanged, while the historical check and its failed observations
lead to a revision satisfying the destination requirement. Detailed inputs,
source identities, and interventions appear in Appendix~\ref{sec:development-details}.

\paragraph{Content changes which behavior is recovered.}
\label{sec:content}
For the constraint and timezone cases, four fresh complete agents per task
receive no memory, a requirement, an implementation reference, or both.
The source and delivery are fixed before editing. All eight repairs pass
official verification (Table~\ref{tab:memorycontent}). On the constraint task,
the requirement-only agent constructs a nearby passing example on different
columns, losing the same-column trigger. On the timezone task, all memory
conditions recover the destination requirement, but implementation-bearing
inputs introduce different errors on later MySQL and named-zone diagnostics.
Those diagnostics were constructed after inspecting the repairs and applied
to all four submissions; they are not predeclared population-level outcomes.

\begin{table}[!htbp]
\centering
\caption{Memory-content interventions: one fresh attempt per input and development task.
All eight pass official verification. MySQL and named-zone checks were constructed
after inspecting the timezone repairs and applied to all four submissions.}
\label{tab:memorycontent}
\small
\begin{tabular}{lcccc}
\toprule
Memory input & Constraint joint & Timezone joint & MySQL cases & Named-zone cases \\
\midrule
None & Fail & Fail & Pass & Fail \\
Requirement & Fail & Pass & Pass & Pass \\
Implementation & Pass & Pass & Pass & Fail \\
Both & Pass & Pass & Fail & Pass \\
\bottomrule
\end{tabular}
\end{table}

A sentence can nevertheless suffice. On the redirect task, requirement-only
and requirement-plus-feedback refinements find the same correct implementation.
Two subsequent original-base constraint agents, with and without a
source-validated example, also pass the joint checks. The role of executable
memory is to preserve and test a decisive condition when adaptation loses it,
rather than to make every memory payload maximally long.

\paragraph{Supplying feedback recovers a specific missed condition.}
Two fresh agents refine the same failed constraint Requirement candidate,
both receiving the requirement again. The augmented agent also receives the
unchanged source check and actual failures for a primary-key or unique field.
Requirement-only review excludes the primary key but still fails the ordinary
unique-field condition. Feedback adds generated-name disambiguation and passes
both. The revisions both pass official verification and the withheld joint
sequence. The observed difference is recovery of the condition identified by
feedback, with the starting candidate and historical requirement held fixed.

\paragraph{Relations preserve the conditions behind the values.}
\label{sec:composition-results}
On Django12663, two retrieval packets both contain source13590, yet their
initial repairs miss the named-tuple behavior. Revisions from the same candidate
show why content matters: execution feedback recovers named tuples but breaks
ordinary tuple/list subclasses; original history preserves the constructor
distinction and all displayed behavior groups
(Appendix Figure~\ref{fig:memory-content-scope}). Retrieving the right source
therefore does not guarantee that the agent uses its condition correctly.

A history-informed Django12663 candidate passes separate
current and historical examples but fails their combination. Four fresh
refinements recover all seven measured conditions in one of two history-only
attempts and two of two attempts given the joint program and its failure.
All four pass official verification. The deterministic range adapter then
reconstructs the same omitted interaction while retaining the lazy User and
nested annotation; graph delivery recovers this program from the source relation.

\paragraph{A matched pair isolates added conditions and observations.}
An original-base pair receives identical history, public scalar program,
observed failure, execution tool, and instructions. Only the augmented arm
also receives the composed programs and their failures
(Table~\ref{tab:composedconditions}). Both repairs pass official verification.
The augmented repair also passes both historical examples and all withheld
joint queries while preserving ordinary subclasses. This comparison tests
delivery of composed checks plus feedback; \expmark{the shared-composer relation contrast C--B, composer contrast C--G, and
feedback contrast F--C remain pending in Table~\ref{tab:planned-graph}.}

\begin{table}[!htbp]
\centering
\caption{Added conditions in the original-base Django12663 pair. Both inputs
contain the same history, current program and failure, tool, and execution
instruction. Only the second adds the composed programs and their failures.
Each row counts the named checks, not independent tasks.}
\label{tab:composedconditions}
\begin{tabular}{lcc}
\toprule
Measured behavior & Current feedback & + Composed conditions \\
\midrule
Official target verification & Pass & Pass \\
Historical named-tuple examples & 0/2 & 2/2 \\
Withheld nonempty joint queries & 2/3 & 3/3 \\
Withheld ordinary subclasses & 2/2 & 2/2 \\
Mechanical scalar / named tuple / tuple & 2/3 & 3/3 \\
\bottomrule
\end{tabular}
\end{table}

The same case also explains the retrieval design: an early hard operation
match rejects source13590 even though FAISS ranks it first. We therefore retain
similarity candidates and use graph relations for adaptation rather than as
an unconditional retrieval veto. A second role-composition example varies
database, ambient, and destination timezones across 72 source conditions.
Requirement and Both pass all 72; Implementation fails 32, including 16
passed without memory. Appendix~\ref{sec:development-details} preserves the
full construction and its source before/after outcomes.

\subsection{Experiment 3: generalization, scale, and efficiency}
\label{sec:recorded-generalization}
\paragraph{Transfer across backbones.}
The historical Related-Lite98 single-shot protocol compares six conditions
using the same 300-summary pool within each backbone
(Table~\ref{tab:swectx-crossmodel}). ContextGraph has the highest resolution
rate among its non-oracle memory conditions on all four models. GPT-5.4
resolves 34/98 versus 21/98 without memory. Its 34/98 is close to the
oracle-summary condition's 35/98, but ContextGraph can retrieve several
summaries whereas the oracle supplies one paired summary; this is a system
comparison rather than a retrieval-accuracy estimate.

\begin{table}[!htbp]
\centering
\caption{Historical Related-Lite98, one attempt per task ($V_1$).
Cells show resolved/completed; bold marks the best non-oracle condition by rate.
Memory conditions draw from the same 300-summary pool within each backbone.}
\label{tab:swectx-crossmodel}
\begin{tabular}{lcccc}
\toprule
Memory & GPT-5.4 & DeepSeek V4 Pro & Kimi K2.6 & MiniMax M2.7 \\
\midrule
None & 21/98 & 30/98 & 18/97 & 19/98 \\
Random summaries & 19/97 & 26/96 & 25/98 & 19/97 \\
Mem0 & 19/98 & 26/97 & 25/98 & 18/98 \\
LangMem & 24/98 & 33/97 & 24/98 & 21/98 \\
Oracle summary & 35/98 & 27/97 & 24/98 & 15/95 \\
ContextGraph & \textbf{34/98} & \textbf{34/97} & \textbf{28/98} & \textbf{27/98} \\
\bottomrule
\end{tabular}
\end{table}

\paragraph{Earlier resolution under retries.}
A separate Related-Lite99 campaign uses DeepSeek V4 Pro with 500 calls and
\$30 per task in both arms. Its 291 curated static episodes provide reference
fixes, verification hints, and pre-base history before approach selection and
after an error. Each later pass retries only that arm's unresolved tasks.
Memory reaches 25 cumulative resolutions in two passes, while the control
requires three (Table~\ref{tab:swectx-stream}). The first two passes yield seven
and nine additional resolutions, respectively; at pass three the gap is two.
Excluding ten targets with near-duplicate source fixes retains the first-pass
gain, 19/89 versus 13/89 ($p=.031$). Additional backbone replications appear
in Appendix Table~\ref{tab:swectx-xmodel}.

\begin{table}[!htbp]
\centering
\caption{Historical Related-Lite99 with curated episodic memory.
Resolved counts are cumulative. Wins/losses compare ContextGraph with control;
$p$ is the paired exact McNemar test.}
\label{tab:swectx-stream}
\begin{tabular}{lcccc}
\toprule
Retry budget & No memory & ContextGraph & Wins/losses & $p$ \\
\midrule
1 pass & 14/99 & \textbf{21/99} & 7/0 & .016 \\
2 passes & 16/99 & \textbf{25/99} & 9/0 & .0039 \\
3 passes & 25/99 & 27/99 & 6/4 & .75 \\
\bottomrule
\end{tabular}
\end{table}

\paragraph{Source availability and interference.}
With the historical GPT-5.4 agent and retrieval procedure fixed, the full
300-summary pool yields the best observed result, 34/98
(Table~\ref{tab:swectx-poolsize}). Smaller pools yield 21--27 resolutions.
Growing a pool changes both source coverage and ranking competition;
\expmark{Table~\ref{tab:planned-pool} separates distractor count from useful-source
availability, and Table~\ref{tab:planned-retrieval} compares dense, lexical,
fused, and relation-aware ranking under fixed $k$ and delivery format.}

\begin{table}[!b]
\centering
\caption{Historical Related-Lite98 memory-pool ablation, GPT-5.4, single-shot
$V_1$. The separate best-of-two $V_2$ result is in Appendix Table~\ref{tab:swectx-v2}.}
\label{tab:swectx-poolsize}
\begin{tabular}{lrrrr}
\toprule
Number of summaries & 50 & 100 & 200 & 300 \\
\midrule
Resolved/completed & 21/96 & 27/98 & 26/97 & \textbf{34/98} \\
Resolution rate & 21.88\% & 27.55\% & 26.80\% & \textbf{34.69\%} \\
\bottomrule
\end{tabular}
\end{table}

\paragraph{Cost and robustness to be measured under the same protocol.}
Fewer retry rounds do not establish lower total cost. \expmark{The planned efficiency
and amortization tables count build, selection, condition execution, and failed
attempts, and distinguish raw tokens, charged USD, and active time.
The second-backbone and sensitivity panels repeat the matched methods and
report outcomes together with coverage.} Appendix~\ref{sec:new-results} retains
the other supplied cross-model and token summaries with their missing settings.
